% GENERATED FILE. Edit canonical lessons, then run npm run generate:books.
% canonical-source-hash: fnv1a64:d00982e3b7583cde
% canonical-lessons: TE-C186-manavadu, TE-C186-manavaralu, TE-C186-tallidandrulu, TE-C186-kavalalu, TE-C186-papa

\chapter{Grandson, Granddaughter, Parents, Twins, Baby}
\label{ch:te-grandson-granddaughter-parents-twins-baby}
\hlchaptermodality{186}

\begin{chapteropening}
\textbf{By the end of this chapter:} \emph{I can say a grandson, a granddaughter, parents, twins and a baby.}

Complete the last lesson of chapter 186: I can say a grandson, a granddaughter, parents, twins and a baby.
\end{chapteropening}

\section[\texorpdfstring{\te{మనవడు} (manavaḍu)}{మనవడు (manavaḍu)}]{\te{మనవడు} (manavaḍu) — a grandson}

\label{lesson:TE-C186-manavadu}



\begin{quote}
Before the new one: say the Telugu for the post, then the Telugu for a pencil.
\end{quote}

\subsection*{You'll want to know: \te{మనవడు}}
\textbf{\te{మనవడు}} — \emph{manavaḍu} — \textquotedblleft{}a grandson\textquotedblright{}.

\begin{grammarlens}[title={What you've built}]
One more word for this chapter.
\end{grammarlens}

\subsection*{Your turn}
\begin{itemize}
\raggedright
  \item \emph{Say it:} \emph{manavaḍu}
  \item \emph{Say it:} \emph{manavaḍu}, once more
  \item \emph{Say it:} say \emph{manavaḍu}
  \item \emph{Recall:} say the Telugu for the post, then the Telugu for a pencil, then say \emph{manavaḍu} again
\end{itemize}

\begin{tcolorbox}[breakable,colback=teal!4,colframe=teal!35!black,title={Before you move on}]
What does \te{మనవడు} mean? (\textquotedblleft{}A grandson\textquotedblright{}.) Say it once more.
\end{tcolorbox}
\section[\texorpdfstring{\te{మనవరాలు} (manavarālu)}{మనవరాలు (manavarālu)}]{\te{మనవరాలు} (manavarālu) — a granddaughter}

\label{lesson:TE-C186-manavaralu}



\begin{quote}
Before the new one: say the Telugu for a pencil, then the Telugu for a grandson.
\end{quote}

\subsection*{You'll want to know: \te{మనవరాలు}}
\textbf{\te{మనవరాలు}} — \emph{manavarālu} — \textquotedblleft{}a granddaughter\textquotedblright{}.

\begin{grammarlens}[title={What you've built}]
One more word for this chapter.
\end{grammarlens}

\subsection*{Your turn}
\begin{itemize}
\raggedright
  \item \emph{Say it:} \emph{manavarālu}
  \item \emph{Say it:} \emph{manavarālu}, once more
  \item \emph{Say it:} say \emph{manavarālu}
  \item \emph{Recall:} say the Telugu for a pencil, then the Telugu for a grandson, then say \emph{manavarālu} again
\end{itemize}

\begin{tcolorbox}[breakable,colback=teal!4,colframe=teal!35!black,title={Before you move on}]
What does \te{మనవరాలు} mean? (\textquotedblleft{}A granddaughter\textquotedblright{}.) Say it once more.
\end{tcolorbox}
\section[\texorpdfstring{\te{తల్లిదండ్రులు} (tallidaṇḍrulu)}{తల్లిదండ్రులు (tallidaṇḍrulu)}]{\te{తల్లిదండ్రులు} (tallidaṇḍrulu) — parents}

\label{lesson:TE-C186-tallidandrulu}



\begin{quote}
Before the new one: say the Telugu for a grandson, then the Telugu for a granddaughter.
\end{quote}

\subsection*{You'll want to know: \te{తల్లిదండ్రులు}}
\textbf{\te{తల్లిదండ్రులు}} — \emph{tallidaṇḍrulu} — \textquotedblleft{}parents\textquotedblright{}.

\begin{grammarlens}[title={What you've built}]
One more word for this chapter.
\end{grammarlens}

\subsection*{Your turn}
\begin{itemize}
\raggedright
  \item \emph{Say it:} \emph{tallidaṇḍrulu}
  \item \emph{Say it:} \emph{tallidaṇḍrulu}, once more
  \item \emph{Say it:} say \emph{tallidaṇḍrulu}
  \item \emph{Recall:} say the Telugu for a grandson, then the Telugu for a granddaughter, then say \emph{tallidaṇḍrulu} again
\end{itemize}

\begin{tcolorbox}[breakable,colback=teal!4,colframe=teal!35!black,title={Before you move on}]
What does \te{తల్లిదండ్రులు} mean? (\textquotedblleft{}Parents\textquotedblright{}.) Say it once more.
\end{tcolorbox}
\section[\texorpdfstring{\te{కవలలు} (kavalalu)}{కవలలు (kavalalu)}]{\te{కవలలు} (kavalalu) — twins}

\label{lesson:TE-C186-kavalalu}



\begin{quote}
Before the new one: say the Telugu for a granddaughter, then the Telugu for parents.
\end{quote}

\subsection*{You'll want to know: \te{కవలలు}}
\textbf{\te{కవలలు}} — \emph{kavalalu} — \textquotedblleft{}twins\textquotedblright{}.

\begin{grammarlens}[title={What you've built}]
One more word for this chapter.
\end{grammarlens}

\subsection*{Your turn}
\begin{itemize}
\raggedright
  \item \emph{Say it:} \emph{kavalalu}
  \item \emph{Say it:} \emph{kavalalu}, once more
  \item \emph{Say it:} say \emph{kavalalu}
  \item \emph{Recall:} say the Telugu for a granddaughter, then the Telugu for parents, then say \emph{kavalalu} again
\end{itemize}

\begin{tcolorbox}[breakable,colback=teal!4,colframe=teal!35!black,title={Before you move on}]
What does \te{కవలలు} mean? (\textquotedblleft{}Twins\textquotedblright{}.) Say it once more.
\end{tcolorbox}
\section[\texorpdfstring{\te{పాప} (pāpa)}{పాప (pāpa)}]{\te{పాప} (pāpa) — a baby, a little one}

\label{lesson:TE-C186-papa}



\begin{quote}
Before the new one: say the Telugu for parents, then the Telugu for twins.
\end{quote}

\subsection*{You'll want to know: \te{పాప}}
\textbf{\te{పాప}} — \emph{pāpa} — \textquotedblleft{}a baby, a little one\textquotedblright{}.

\begin{grammarlens}[title={What you've built}]
One more word for this chapter.
\end{grammarlens}

\subsection*{Your turn}
\begin{itemize}
\raggedright
  \item \emph{Say it:} \emph{pāpa}
  \item \emph{Say it:} \emph{pāpa}, once more
  \item \emph{Say it:} say \emph{pāpa}
  \item \emph{Recall:} say the Telugu for parents, then the Telugu for twins, then say \emph{pāpa} again
\end{itemize}

\begin{tcolorbox}[breakable,colback=teal!4,colframe=teal!35!black,title={Before you move on}]
What does \te{పాప} mean? (\textquotedblleft{}A baby, a little one\textquotedblright{}.) Say it once more.
\end{tcolorbox}
